\documentclass[yuxi]{../template/Report}%方括号内写yuxi即生成预习报告\documentclass[yuxi]{../template/Report}
\settemplatedir{../template/}%设置模板路径
\exname{双棱镜干涉} %实验名称
\extable{} %实验桌号
\instructor{殷立明} %指导教师
\class{微电子2401} %班级
\name{黄千格} %姓名
\stuid{3240100650} %学号

\nyear{2025} %年
\nmonth{12} %月
\nday{22} %日
\nweekday{一} %星期几，e.g. \nweekday{三}
\daypart{上午}%上午/下午

\redate{} %如有实验补做，补做日期
\resitu{} %情况说明：
\begin{document}
\maketitle

\section{预习报告（10分）}

\subsection{实验综述}

\subsubsection{实验原理}

\paragraph{1.光的干涉条件与方法}
光的干涉需要满足三个基本条件：
\begin{enumerate}
    \item \textbf{频率相同} ($f_1 = f_2$)；
    \item \textbf{振动方向平行} ($\vec{A}_1 // \vec{A}_2$)；
    \item \textbf{相位差恒定} ($\phi_1 - \phi_2 = C$)。
\end{enumerate}

获得相干光的方法主要有两种：
\begin{itemize}
    \item \textbf{分波前法}：如杨氏双缝干涉、菲涅尔双棱镜干涉（本实验采用）。
    \item \textbf{分振幅法}：如薄膜干涉、迈克尔逊干涉仪。
\end{itemize}

\begin{figure}[H]
    \centering
    \begin{minipage}[b]{0.48\textwidth}
        \centering
        \includegraphics[width=\linewidth]{figures/1.jpg}
        \caption{分波前法}
    \end{minipage}\hfill
    \begin{minipage}[b]{0.48\textwidth}
        \centering
        \includegraphics[width=\linewidth]{figures/2.jpg}
        \caption{分振幅法}
    \end{minipage}

\end{figure}

\paragraph{2.明暗条纹条件}
根据光程差 $\delta$，干涉条纹的分布满足以下条件：
\begin{equation}
    \delta = 
    \begin{cases} 
    \pm 2k \cdot \frac{\lambda}{2}, & k = 0, 1, 2 \cdots \quad (\text{明纹}) \\
    \pm (2k + 1) \frac{\lambda}{2}, & k = 1, 2 \cdots \quad (\text{暗纹})
    \end{cases}
\end{equation}

\paragraph{3.波长测量原理}
菲涅尔双棱镜将一个光源 $S$ 分为两个虚光源 $S_1$ 和 $S_2$。
设虚光源间距为 $d$，虚光源到屏幕（或焦平面）的距离为 $D$，相邻明纹（或暗纹）的间距为 $\Delta x$，则波长 $\lambda$ 可由下式计算：
\begin{equation}
    \Delta x = \frac{\lambda D}{d} \quad \Rightarrow \quad \lambda = \frac{\Delta x \cdot d}{D}
\end{equation}
其中，由于虚光源 $S_1, S_2$ 是不可见的，直接测量 $d$ 和 $D$ 较为困难，因此实验中采用二次成像法（贝塞尔法）进行测量。

\paragraph{4.二次成像法测 $d$ 和 $D$}
在双棱镜与屏幕之间加入一个凸透镜，移动透镜位置，可以在屏幕上分别得到一次放大的像和一次缩小的像。
\begin{itemize}
    \item 设大像的间距为 $d_1$，小像的间距为 $d_2$，则虚光源间距 $d$ 为：
    \begin{equation}
        d = \sqrt{d_1 d_2}
    \end{equation}
    \item 设透镜两次成像时，物像距离之和（即狭缝到毛玻璃屏的距离）为 $D$。在实验装置中，通常 $D = v_1 + v_2$（即狭缝到屏幕的距离）。
    \item 最终波长计算公式为：
    \begin{equation}
        \lambda = \frac{\Delta x \cdot \sqrt{d_1 d_2}}{D}
    \end{equation}
\end{itemize}
\textbf{注意}：为了能成两次像，必须满足条件 $D > 4f$。

\begin{figure}[H]
    \centering
    % 占位符：插入PPT中二次成像法的光路图
    \includegraphics[width=0.7\textwidth]{figures/3.jpg} 
    \caption{二次成像法测量光路原理图}
\end{figure}

\subsubsection{实验内容}
\paragraph{0.预备实验}

\paragraph{(1) 透镜位置调节（共轴调节）}
在正式测量前，需要确保透镜光轴与激光束重合。
\begin{enumerate}
    \item 放上透镜，在后方放置带有小孔的白屏。
    \item 调节透镜支架上的\textbf{1号螺丝}（可调高低）和\textbf{2号旋钮}（可调左右）。
    \item \textbf{调节标准}：使得激光光点经过透镜折射后，仍然能通过白屏上的定位小孔，确光路共轴。
\end{enumerate}


\paragraph{(2)测量凸透镜焦距（位移法/贝塞尔法）}
利用位移法测量透镜焦距，验证透镜成像规律，为后续二次成像测 $d$ 做准备。
\begin{enumerate}
    \item 原理公式：
    根据凸透镜成像公式 $\frac{1}{f} = \frac{1}{u} + \frac{1}{v}$。固定物与屏之间的距离 $D$（需满足 $D > 4f$），移动透镜，可在屏上找到两次清晰成像的位置。
    设两次成像时透镜移动的距离为 $\Delta l$，则焦距 $f$ 为：
    \begin{equation}
        f = \frac{D^2 - \Delta l^2}{4D}
    \end{equation}
    \item 操作步骤：
    \begin{itemize}
        \item 固定物（光源）和像屏，测量间距 $D$。
        \item 移动透镜，找到第一个清晰像（放大像），记下位置。
        \item 继续移动透镜，找到第二个清晰像（缩小像），记下位置。
        \item 计算透镜位移量 $\Delta l$，代入公式求 $f$。
    \end{itemize}
    \item 注意：$D$ 越大，两次成像位置差异越大；若 $D=4f$，仅成一次像。
\end{enumerate}

\paragraph{1. 光路粗调}
\begin{enumerate}
    \item \textbf{狭缝调节}：将毛玻璃屏置于靠近摄像头位置，狭缝到屏幕距离取 $126.00\,cm$。调节螺丝使激光通过狭缝，并改变狭缝大小以出现衍射条纹。调节狭缝方向使衍射条纹水平（即狭缝竖直）。
    \item \textbf{双棱镜位置}：将双棱镜置于距离狭缝 $22.00\,cm$ 处。调节双棱镜使激光垂直入射，并调节棱脊竖直（此时干涉条纹竖直）。
    \item \textbf{干涉区对中}：调节双棱镜的左右位置，使干涉区域位于光场中间，确保二次成像时两个像点的明暗程度一致。
    \item \textbf{摄像头调整}：调节摄像头高低、左右位置，使其正对干涉区域，旋转镜头使叉丝竖线与条纹平行。
\end{enumerate}


\paragraph{2. 二次成像法测 $D$ 和 $d$}
\begin{enumerate}
    \item 放入凸透镜，在双棱镜和毛玻璃屏之间水平移动。
    \item 观察屏幕，找到像点最清晰的两个透镜位置（一处像间距大，一处像间距小）。
    \item 移动摄像头，利用叉丝竖线对准像点\textbf{中间位置}，分别测量 $d_1$ 和 $d_2$。
    \item 记录此时狭缝到屏幕的距离 $D$（或通过透镜位置计算，具体依仪器读数方式而定）。
\end{enumerate}



\paragraph{3. 测定干涉条纹间距 $\Delta x$}
\begin{enumerate}
    \item 移开透镜，观察干涉条纹。
    \item 若条纹不清或太暗，可缓慢微调狭缝宽度（太宽会导致条纹模糊，太窄会导致亮度不足）。
    \item 使用摄像头底座的螺旋测微器测量条纹位置。
    \item \textbf{测量方法}：转动鼓轮，使叉丝竖线对齐条纹位置。为了提高精度，应针对\textbf{暗纹两边的明暗交界线}各读一个数取平均值作为该级条纹的位置。
    \item 连续测量多个条纹的位置，利用逐差法或线性回归求出 $\Delta x$。
\end{enumerate}

\subsubsection{实验现象}
\begin{enumerate}
    \item \textbf{狭缝衍射}：仅放置狭缝时，在屏幕上观察到激光的单缝衍射图样。
    \item \textbf{虚光源成像}：加入凸透镜进行二次成像时，屏幕上会出现两个清晰的红色亮斑（虚光源的像）。一次间距较宽，一次间距较窄。
    \item \textbf{双棱镜干涉}：移去透镜后，在重叠区域观察到明暗相间、等间距的竖直干涉条纹。
\end{enumerate}

\begin{figure}[H]
    \centering
    % 占位符：插入PPT中最终的干涉条纹图片
    \includegraphics[width=0.5\textwidth]{figures/4.jpg} 
    \caption{观察到的双棱镜干涉条纹}
\end{figure}

\section{实验重点}
\begin{enumerate}
    \item \textbf{掌握分波前法干涉原理}：理解菲涅尔双棱镜如何将一束光分为两束相干光。
    \item \textbf{光路调节技巧}：掌握共轴调节的方法，确保激光、狭缝、双棱镜棱脊、透镜光心和观测屏中心在同一高度和轴线上。
    \item \textbf{二次成像法的应用}：理解并掌握如何利用辅助透镜测量不可直接测量的虚光源间距 $d$。
\end{enumerate}

\section{实验难点}
\begin{enumerate}
    \item \textbf{共轴等高调节}：若光路未调节至共轴等高，干涉条纹可能倾斜、不清晰，或者在移动透镜时像点跑出视场，这是实验操作中最耗时且关键的一步。
    \item \textbf{清晰像的判断}：在二次成像法中，准确判断像点“最清晰”的位置存在主观误差，需要仔细调节。
    \item \textbf{条纹间距的精确测量}：由于条纹具有一定的宽度，且存在边缘效应，如何利用叉丝准确对准（对准明暗交界处取平均）以减小读数误差是测量精度的关键。
    \item \textbf{回程误差的消除}：在使用螺旋测微器进行测量时，需单方向旋转鼓轮以消除螺距空程带来的误差。
\end{enumerate}

\begin{fullreportonly}
\newpage


\section{原始数据（20分）}
\begin{figure}[H]
    \centering
    \includegraphics[width=0.55\textwidth]{实验数据原图.jpg}
    \caption{实验数据原图}
    \label{fig:原始数据1}
\end{figure}


\section{结果与分析（60分）}
\subsection{数据处理与结果（30分）}


\subsection{误差分析（20分）}

\section{思考题（10分）}

\end{fullreportonly}
\insertnotes
\end{document}
